\section{Introduction}

Language-model agents are moving from stateless responders to experience-driven
systems that improve over time by distilling their own interaction trajectories
into a persistent memory bank and reusing it through a retrieve--execute--write
loop~\tbd{CITATION_REASONINGBANK}. The implicit premise of this paradigm is that
the system knows \emph{which memory deserves credit}: writing, merging, evicting,
and retrieving all depend on a per-memory utility signal. Yet the signals in
widely used systems are, on the characterization we adopt from prior descriptions
and subject to bibliographic verification, \emph{correlational}. When a task
succeeds and a memory was retrieved, that memory's utility is incremented. This
conflates three distinct events: a memory can be \emph{retrieved} without being
\emph{used}, and used without providing genuine \emph{causal} help.

\paragraph{Retrieval is not contribution.}
A memory's \emph{immediate conditional retrieval contribution} is defined only
counterfactually---by how much the immediate task outcome degrades when the memory
is removed from the retrieved set, conditioned on the current query, retrieved set,
and deployment history. Correlational credit does not measure this quantity. We
hypothesize that it behaves like the reinforcement
schedule in Skinner's superstition experiment~\tbd{CITATION_SKINNER_SUPERSTITION}:
memories that merely co-occur with success are persistently strengthened, so that
\emph{superstitious memories}---entries with zero or negative causal
contribution that nonetheless rise into the top-utility set and occupy retrieval
slots---are not corrected. Because retrieval probability tends to track query
solvability and high utility feeds back into higher retrieval rank, we hypothesize a
self-reinforcing loop in which the fraction of superstitious memories \emph{grows}
with deployment length. A second, related failure is cross-task interference: a
memory that is genuinely useful for one task family is indiscriminately retrieved
into another and harms the new task.

\paragraph{Why prior signals are insufficient.}
Existing approaches in our comparison commonly use observational, co-occurrence
credit: success-trajectory distillation~\tbd{CITATION_REASONINGBANK}, environment-feedback
utility~\tbd{CITATION_SKEMEX}, and Monte-Carlo or reinforcement updates over
outcomes~\tbd{CITATION_MEMRL}. What is learned is the management \emph{policy};
the underlying signal remains correlational. We distinguish passive co-occurrence
counting from outcome-reward reinforcement learning under exploration, which is
correlational only in a weaker sense; the diagnostic below is scoped to the systems
and streams we evaluate, not to all learned managers. Counterfactual data-valuation
methods such as Data Shapley~\tbd{CITATION_DATA_SHAPLEY}, influence
functions~\tbd{CITATION_INFLUENCE_FUNCTIONS}, and TracIn~\tbd{CITATION_TRACIN}
target static, offline training data, while context-attribution
methods~\tbd{CITATION_CONTEXTCITE} explain a single generation. None, to the best
of our current and still-to-be-verified survey, deliver counterfactual credit in
an online, non-stationary, closed loop where attribution itself changes the data
distribution.

\paragraph{Approach.}
We propose \method{}, which replaces correlational bookkeeping with counterfactual
credit while keeping it affordable. \method{} adds three tightly coupled
components to the standard loop, in a strictly linear dependency
\rit{}~$\rightarrow$~\aca{}~$\rightarrow$~consumers: (i) \emph{Randomized
Interventional Trials} (\rit{}), a small-budget paired leave-one-out protocol that
produces interventional outcome-difference labels; (ii) \emph{Amortized Counterfactual Attribution} (\aca{}),
a per-event attributor whose scoring cost is constant in bank size (at fixed $d$
and top-$k$), trained on those labels and periodically
recalibrated to track closed-loop drift; and (iii) two deliberately simple
consumers of the resulting score---threshold-based credit governance and
scope-gated retrieval. Figure~\ref{fig:problem-solution} contrasts correlational
credit with this pipeline. Our central position is that reliable \emph{immediate
counterfactual credit}---not better heuristics---is a missing primitive for the
credit-assignment step of self-evolution, in the immediate-contribution setting we
evaluate; governance and scope gating are then empirically tested consumers of that
signal, not validated causal estimators of long-horizon retention value.

\paragraph{Contributions.}
\begin{itemize}
  \item \textbf{(C1) Superstitious memories.} We formalize and plan to quantify a
  diagnostic phenomenon in the memory systems and streams we evaluate: under
  correlational credit, zero/negative-contribution
  memories accumulate in the top-utility set and their fraction is hypothesized to
  rise with deployment time, and
  we introduce an interventional audit (\rit{}) with two statistics---the
  credit--contribution correlation (CCC) and the superstition rate (SR@$k$)---to
  measure it.
  \item \textbf{(C2) Amortized counterfactual attribution.} We introduce \aca{},
  which amortizes interventional labels into a per-event contribution
  score that is constant in bank size (at fixed feature dimension $d$ and top-$k$),
  with periodic recalibration, and we plan to test whether it recovers
  causal credit substantially better than the correlational baseline at low cost.
  \item \textbf{(C3) Directly consumable causal credit.} We contribute a design in
  which a single causal-credit signal is directly consumable: threshold governance
  and scope-gated retrieval are two consumers of the same score. We plan to test
  whether they respectively flatten superstition accumulation and reduce
  cross-task interference. Substituting the score for a manager's outcome reward is
  reported only as a secondary signal-portability check, not as a separate
  contribution.
  \item \textbf{(C4) Systematic evaluation.} We plan a systematic evaluation
  across four benchmark categories, three backbones, and three seeds, testing
  whether \method{} improves average success over the strongest baseline at low
  token overhead.
\end{itemize}

\paragraph{Planned evidence.}
Every empirical statement in this paper is a planned test, not a measured result;
unresolved quantities are written as registered ledger placeholders. A diagnostic
pilot (Section~4) tests whether superstition accumulates; the mechanism analysis
(Section~6) tests whether \aca{} recovers the audited signal in closed loop under
the same axes as the pilot; and the main study tests whether the recovered signal
translates into a planned average gain of \tbd{MAIN_AVG_GAIN_OVER_BEST_BASELINE}
over the strongest baseline. Each claim carries an explicit falsification gate,
and results that fail a gate will narrow the corresponding claim rather than be
suppressed.

% Ledger IDs: [[TBD:FIG1_TOPUTILITY_ZERO_EFFECT_FRACTION]]
\begin{figure*}[t]
\centering
\figureplaceholder{1.35in}{%
\textbf{Fig. 1 placeholder: Problem and solution schematic}\\[3pt]
Left panel: correlational credit and a top-utility memory set, with memories of non-positive causal contribution highlighted.\\
Right panel: the \rit{} $\rightarrow$ \aca{} $\rightarrow$ consumers pipeline on the retrieve-execute-write loop.\\
Future source: pilot SR@k and mechanism analysis. No result geometry is drawn at this stage.}
\caption{Planned hero schematic contrasting correlational credit with the counterfactual-credit pipeline. The eventual empirical annotation remains TBD until verification.}
\label{fig:problem-solution}
\end{figure*}

